% Generated draft: requires mathematical and source review.
\documentclass[11pt]{article}
\usepackage[margin=25mm]{geometry}
\usepackage{fontspec}
\setmainfont{Latin Modern Roman}
\usepackage{amsmath,amssymb,mathtools,unicode-math}
\setmathfont{Latin Modern Math}
\usepackage{enumitem,needspace,xurl,hyperref}
\setlength{\parindent}{0pt}
\setlength{\parskip}{5pt}
\setlength{\emergencystretch}{3em}
\newcommand{\N}{\mathbb{N}}
\newcommand{\Z}{\mathbb{Z}}
\newcommand{\Q}{\mathbb{Q}}
\newcommand{\R}{\mathbb{R}}
\newcommand{\C}{\mathbb{C}}
\newcommand{\F}{\mathbb{F}}
\newcommand{\A}{\mathbb{A}}
\newcommand{\PP}{\mathbb{P}}
\newcommand{\E}{\mathbb{E}}
\newcommand{\eps}{\varepsilon}
\newcommand{\dd}{\,\mathrm{d}}
\begin{document}
\section*{UnsolvedMath ID: 30005498}
\textbf{Source problem code:} catalogue entry
\subsection{Definitions and mathematical statement}
The complex-analytic language of the supplied excerpt indicates the following intended setting, although the base field is not explicitly specified in the problem statement. Let \((A,\Theta)\) be a complex principally polarized abelian variety, and put \(g=\dim_{\C}A\). Here \(\Theta\) is an effective ample divisor whose associated polarization \(A\to\operatorname{Pic}^{0}(A)\) is an isomorphism. The polarized variety is \emph{indecomposable} if it is not isomorphic, as a polarized abelian variety, to a product of two positive-dimensional principally polarized abelian varieties.

For an integer \(m\geq2\), define the multiplicity locus
\[
\operatorname{Sing}_{m}(\Theta)
 :=\{x\in A:\operatorname{mult}_{x}(\Theta)\geq m\}.
\]
The multiplicity is the order of vanishing of a local defining equation of \(\Theta\) at \(x\). We take the dimension of this locus with its reduced structure and use \(\dim\varnothing=-\infty\). These are the conventional interpretations used here; the supplied excerpt does not explicitly define its multiplicity-locus notation or its empty-set convention.

\textbf{Question.} Is it true that, for every indecomposable complex principally polarized abelian variety \((A,\Theta)\) and every integer \(m\geq2\),
\[
\dim\operatorname{Sing}_{m}(\Theta)\leq g-2m+1\,?
\]
In particular, when \(g-2m+1<0\), this asks that the multiplicity locus be empty.

\textbf{Supplied background and qualifications.} The excerpt attributes to Koll\'ar, reference [9, Theorem 17.13] in the report, the statement that the log canonical threshold of \(\Theta\) is \(1\), equivalently that \((A,\Theta)\) is log canonical, and the resulting bound \(\dim\operatorname{Sing}_{m}(\Theta)\leq g-m\) for all integers \(m\geq2\). It attributes to Ein and Lazarsfeld [7], under indecomposability, normality and rational singularities of \(\Theta\), together with the bound \(\dim\operatorname{Sing}_{m}(\Theta)\leq g-m-1\). It attributes to Saito [15] the equivalence between rational singularities of \(\Theta\) and \(\widetilde\alpha_{\Theta}>1\), where \(\widetilde\alpha_{\Theta}\) denotes its minimal exponent. These reference labels are internal to the report; their full bibliographic entries are not supplied.

The excerpt reports that the proposed stronger bound holds for Prym varieties associated with \'etale double covers, by work of Casalaina-Martin [5], and states it as a corollary for \(g\leq5\), with the supplied justification that every principally polarized abelian variety in those dimensions is a Prym variety. The precise scope of that Prym representation assertion is not explained in the record. It also attributes to Musta\c{t}\u{a} and Popa [13] the proposed bound when \(\Theta\) has only isolated singularities, using Hodge ideals.

The excerpt additionally reports a conditional theorem from joint work with Christian Schnell: if the ``center of minimal exponent'' is a one-dimensional scheme \(Y\), then \(Y\) is a smooth hyperelliptic curve and the proposed bound holds for every integer \(m\geq2\). It further states that equality for some such \(m\) implies either \(A=\operatorname{Jac}(Y)\), or
\[
 g=2m,\qquad g(Y)=2m-1,\qquad
 \widetilde\alpha_{\Theta}=\frac{2m-1}{m},
\]
with \(\Theta\) having a unique singular point of multiplicity \(m\). This is retained as a supplied source claim, with unresolved qualifications: the excerpt breaks off before defining the center, and it does not specify the sense of the Jacobian identification. Moreover, its second alternative gives \(\dim\operatorname{Sing}_{m}(\Theta)=0\) while \(g-2m+1=1\), so that alternative cannot itself realize equality under the multiplicity-locus interpretation above. The exact conditional equality statement therefore requires checking against the complete source. The general equality classification mentioned in the supplied assessment is not formulated precisely enough here to constitute an additional question.
\subsection{Short English statement}
The question asks whether the points where a theta divisor has multiplicity at least \(m\) always form a locus of dimension at most \(\dim A-2m+1\), for every integer \(m\geq2\), when the principally polarized abelian variety is indecomposable. The formulation uses the complex setting suggested by the excerpt. The supplied record labels the problem partially solved; its special cases, conditional results, and equality discussion remain attributed source claims requiring review.
\subsection{Sources}
\begin{itemize}
\item Supplied local UnsolvedMath registry record, permanent ID \(30005498\), title \emph{Dimension Bounds for Theta Divisor Singularities}, exact source problem code \texttt{catalogue entry}, in the Oberwolfach Reports Open Problems collection. The proposed year is 2023; no proposer is identified. The registry supplies the labels \texttt{partially\_solved} and \texttt{corrected\_verified}, and says its statement was checked against the official report. These are supplied curation claims, not independent verification in this draft.
\item \emph{Komplexe Analysis -- Differential and Algebraic methods in K\"ahler spaces} (2023), Oberwolfach Reports, pp. 965--1030, DOI 10.4171/owr/2023/18. Supplied source link: \url{https://doi.org/10.4171/owr/2023/18}. The supplied excerpt ends at a page marker for p. 1017 and is incomplete.
\item Supplied literature source, labelled ``Schnell and Yang'': \url{https://arxiv.org/abs/2309.16763}. The record supplies no article title or full bibliographic entry.
\item Supplied literature source, labelled ``OWR 2023 report'': \url{https://ems.press/content/serial-article-files/47013}.
\item Literature assessment dated exactly \texttt{2026-08-22}: the record reports that the dimension bound is known for Prym varieties, in dimension at most five, and for isolated singularities, with further conditional results, while the general equality classification remains open. This is supplied dated curation derived from the OWR material, not independently verified; the record expressly requests independent verification of time-sensitive status and cited claims.
\item Background attributions in the supplied excerpt: Koll\'ar [9, Theorem 17.13], Ein and Lazarsfeld [7], Saito [15], Casalaina-Martin [5], and Musta\c{t}\u{a} and Popa [13]. Full entries for these report-internal references are absent. The conditional theorem is described as joint work with Christian Schnell; the speaker is not identified in the excerpt.
\end{itemize}
\end{document}
